% Bibliographic records and theorem locators checked against primary papers,
% author pages, and publisher records on 7 October 2026.
% The bibliography is self-contained: no BibTeX/Biber invocation is needed.
\begin{thebibliography}{99}
\addcontentsline{toc}{section}{References}

\bibitem{abe}
Tetsuya Abe,
\emph{The Rasmussen invariant of a homogeneous knot},
Proceedings of the American Mathematical Society
\textbf{139} (2011), no.~7, 2647--2656.
\href{https://doi.org/10.1090/S0002-9939-2010-10687-1}{doi:10.1090/S0002-9939-2010-10687-1}.
\href{https://arxiv.org/abs/1003.5392}{arXiv:1003.5392}.
Theorems~2.3 and~3.4 give the genus statement and homogeneity criterion used here.

\bibitem{aht}
Ian Agol, Joel Hass, and William P.~Thurston,
\emph{The computational complexity of knot genus and spanning area},
Transactions of the American Mathematical Society
\textbf{358} (2006), no.~9, 3821--3850.
\href{https://doi.org/10.1090/S0002-9947-05-03919-X}{doi:10.1090/S0002-9947-05-03919-X}.
\href{https://arxiv.org/abs/math/0205057v2}{arXiv:math/0205057v2}.
Theorems~12 and~16 describe ordinary and weighted orbit counting.

\bibitem{weaving}
Mark E.~AlSukaiti and Nafaa Chbili,
\emph{Alexander and Jones polynomials of weaving 3-braid links and Whitney
rank polynomials of Lucas lattice},
Heliyon \textbf{10} (2024), no.~7, article e28945.
\href{https://doi.org/10.1016/j.heliyon.2024.e28945}{doi:10.1016/j.heliyon.2024.e28945}.
\href{https://arxiv.org/abs/2303.11398}{arXiv:2303.11398}.
See Corollary~3.5 of the arXiv version for the determinant formula.

\bibitem{bnfast}
Dror Bar-Natan,
\emph{Fast Khovanov homology computations},
Journal of Knot Theory and Its Ramifications
\textbf{16} (2007), no.~3, 243--255.
\href{https://doi.org/10.1142/S0218216507005294}{doi:10.1142/S0218216507005294}.
\href{https://arxiv.org/abs/math/0606318}{arXiv:math/0606318}.
\href{https://www.math.toronto.edu/drorbn/papers/FastKh/FastKh.pdf}{Author's version of 20 May 2007}.

\bibitem{cromwell}
Peter R.~Cromwell,
\emph{Homogeneous links},
Journal of the London Mathematical Society (2)
\textbf{39} (1989), no.~3, 535--552.
\href{https://doi.org/10.1112/jlms/s2-39.3.535}{doi:10.1112/jlms/s2-39.3.535}.

\bibitem{treewidth}
Arnaud de~Mesmay, Jessica Purcell, Saul Schleimer, and Eric Sedgwick,
\emph{On the tree-width of knot diagrams},
Journal of Computational Geometry
\textbf{10} (2019), no.~1, 164--180.
\href{https://doi.org/10.20382/jocg.v10i1a6}{doi:10.20382/jocg.v10i1a6}.
\href{https://arxiv.org/abs/1809.02172v2}{arXiv:1809.02172v2}.
See Theorem~3.7 and Corollary~1.2 for the obstruction to uniformly small
diagrammatic tree-width.

\bibitem{kelomaki}
Tuomas Kelom\"aki and Dirk Sch\"utz,
\emph{On computational complexity of Khovanov homology},
preprint, 5 January 2026, 30~pp.
\href{https://arxiv.org/abs/2601.02119v1}{arXiv:2601.02119v1}.
Theorems~1.2--1.4 distinguish the explicit-basis obstruction,
the specialized three-braid algorithm, and extremal homological bands.

\bibitem{km}
Peter B.~Kronheimer and Tomasz S.~Mrowka,
\emph{Khovanov homology is an unknot-detector},
Publications Math\'ematiques de l'IH\'ES
\textbf{113} (2011), 97--208.
\href{https://doi.org/10.1007/s10240-010-0030-y}{doi:10.1007/s10240-010-0030-y}.
\href{https://arxiv.org/abs/1005.4346}{arXiv:1005.4346}.
The rational-rank implication used here follows from Corollary~1.3,
Proposition~1.4, and the unknot detection discussed in the introduction.

\bibitem{lackenbytalk}
Marc Lackenby,
\emph{Unknot recognition in quasi-polynomial time},
Oxford topology seminar handout, March 2021, 34~pp.
\url{https://people.maths.ox.ac.uk/lackenby/quasipolynomial-talk-oxford-compressed.pdf}.
The announced bound is displayed on slide~7; the four proposed speed-ups
are summarized on slide~25.

\bibitem{lackenby2026}
Marc Lackenby,
\emph{Incompressible surfaces, hierarchies and unknot recognition},
preprint, 25 July 2026, 54~pp.
\href{https://arxiv.org/abs/2607.23350v1}{arXiv:2607.23350v1}.
\href{https://people.maths.ox.ac.uk/lackenby/algorithm-incompressible-250726.pdf}{Author's PDF}.
See Propositions~9.1 and~10.2 for the iteration and hierarchy-length bounds,
and Theorem~14.4 for compressed cutting under uniform-type hypotheses.

\bibitem{lobb}
Andrew Lobb,
\emph{Computable bounds for Rasmussen's concordance invariant},
Compositio Mathematica
\textbf{147} (2011), no.~2, 661--668.
\href{https://doi.org/10.1112/S0010437X10005117}{doi:10.1112/S0010437X10005117}.
\href{https://arxiv.org/abs/0908.2745v2}{arXiv:0908.2745v2}.
The interval used here is Theorem~1.10, with Definitions~1.8--1.9.

\bibitem{manchon}
Pedro M.~G.~Manch\'on,
\emph{Homogeneous links and the Seifert matrix},
Pacific Journal of Mathematics
\textbf{255} (2012), no.~2, 373--392.
\href{https://doi.org/10.2140/pjm.2012.255.373}{doi:10.2140/pjm.2012.255.373}.
\href{https://arxiv.org/abs/1102.0890v2}{arXiv:1102.0890v2}.

\bibitem{morton}
Hugh R.~Morton,
\emph{The multivariable Alexander polynomial for a closed braid},
in \emph{Low-dimensional topology} (Funchal, 1998),
Hanna Nencka, ed., Contemporary Mathematics \textbf{233},
American Mathematical Society, Providence, RI, 1999, 167--172.
\href{https://arxiv.org/abs/math/9803138}{arXiv:math/9803138}.
The introduction and the discussion following Theorem~1 state the
reduced Burau closure formula.

\bibitem{proveit}
Vladimir Reshetnikov,
\emph{ProveIt: Topology/UnknotRecognition},
software and research repository, source snapshot at commit
\pin, accessed 7 October 2026.
\url{https://github.com/VladimirReshetnikov/ProveIt/tree/4e6fe879e5238ec0b134b6d33fdca0b8c7c1c711/Topology/UnknotRecognition}.

\bibitem{shumakovitch}
Alexander N.~Shumakovitch,
\emph{Torsion of Khovanov homology},
Fundamenta Mathematicae \textbf{225} (2014), 343--364.
\href{https://doi.org/10.4064/fm225-1-16}{doi:10.4064/fm225-1-16}.
\href{https://arxiv.org/abs/math/0405474v2}{arXiv:math/0405474v2}
(updated 19 June 2018 to match the published version; arXiv title:
\emph{Torsion of the Khovanov homology}).
Corollary~3.2.C, on published page~356, gives the reduced/unreduced
splitting over $\mathbb F_2$.

\end{thebibliography}
